\documentclass{beamer}


\mode<presentation> {
	

	\usetheme{CambridgeUS}


	
	% As well as themes, the Beamer class has a number of color themes
	% for any slide theme. Uncomment each of these in turn to see how it
	% changes the colors of your current slide theme.
	
	%\usecolortheme{albatross}
	%\usecolortheme{beaver}
	%\usecolortheme{beetle}
	%\usecolortheme{crane}
	%\usecolortheme{dolphin}
	%\usecolortheme{dove}
    %\usecolortheme{fly}
	%\usecolortheme{lily}
	%\usecolortheme{orchid}
	%\usecolortheme{rose}
	%\usecolortheme{seagull}
	%\usecolortheme{seahorse}
	%\usecolortheme{whale}
	%\usecolortheme{wolverine}
	
	%\setbeamertemplate{footline} % To remove the footer line in all slides uncomment this line
	%\setbeamertemplate{footline}[page number] % To replace the footer line in all slides with a simple slide count uncomment this line
	
	%\setbeamertemplate{navigation symbols}{} % To remove the navigation symbols from the bottom of all slides uncomment this line
}


\usepackage{graphicx} % Allows including images
\usepackage{booktabs} % Allows the use of \toprule, \midrule and \bottomrule in tables
\usepackage{amsmath, amsthm, amssymb, caption}
\usepackage{mdframed}
\usepackage{hyperref}
\usepackage{centernot}
\newtheorem{proposition}{Proposition}[section]
\usepackage{multirow}
\setbeamertemplate{theorems}[numbered]
\theoremstyle{plain} % insert bellow all blocks you want in italic
%\newtheorem{theorem}{Theorem}[section] % to number according to section
%\usepackage{enumitem}


\theoremstyle{definition} % insert bellow all blocks you want in normal text
%\newtheorem{definition}{Definition}[section] % to number according to section
\newtheorem*{remark}{Remark} %for remark 
\usepackage[notes,backend=biber]{biblatex-chicago}
%\bibliography{Reference_BGJS}


\AtBeginSection[]
{
  \begin{frame}
    \frametitle{Overview}
    \tableofcontents[currentsection]
  \end{frame}
}




%%%%% Bibliography options %%%%%%%%
\ExecuteBibliographyOptions{giveninits=true, isbn=false, url=false, doi=false, eprint=false, url=false, uniquename=init}
\DeclareBibliographyAlias{webpage}{online}
\AtEveryBibitem{\clearlist{language}}
\AtEveryBibitem{%
      \ifentrytype{online}
        {}
        {\clearfield{urlyear}\clearfield{urlmonth}
\clearfield{doi}	        
\clearfield{urlday}}}


%\hypersetup{%
%  colorlinks=true,
%  linkcolor=blue,
%  linkbordercolor={0 0 1}
%}




%\hypersetup{%
%  colorlinks=true,
%  linkcolor=blue,
%  linkbordercolor={0 0 1}
%}


\setlength{\parindent}{0.0in}
\setlength{\parskip}{0.05in}
%%%%%%%%%%%%%%%%%%%%%%%%%%%%%%%%%%%%%%%%%%%%%%%%%%%%%%%%%% }



%----------------------------------------------------------------------------------------
%	TITLE PAGE
%----------------------------------------------------------------------------------------


\title[ECON 10000 TA Session 3]{Supply and Demand and Consumer Theory} % The short title appears at the bottom of every slide, the full title is only on the title page


\author{} % Your name
\institute[UChicago Econ] % Your institution as it will appear on the bottom of every slide, may be shorthand to save space
{   \normalsize{\href{mailto:}{}}
	% Your email address
}
\date{January 27, 2026} % Date, can be changed to a custom date
\begin{document}


\begin{frame}
	\titlepage % Print the title page as the first slide
\end{frame}

%----------------------------------------------------
\section{Supply and Demand}

\begin{frame}{Review: Aggregate Supply and Demand}
    \textbf{Example:} Concert Tickets vs Spotify Streams

    A Taylor Swift concert has \textbf{50,000} seats (indivisible good), whereas Spotify streams are \textbf{unlimited} (divisible good). What does this tell us about AD and AS?

    For Concerts:
\begin{itemize}
    \item Each consumer buys 0 or 1 ticket
    \item Market Demand = number of people willing to buy at each price
\end{itemize}

    For Spotify:
\begin{itemize}
    \item Each consumer chooses how many hours to stream
    \item Market Demand = horizontal sum of individual quantities
    \item We fix the price and add up how much every consumer wants
\end{itemize}
\end{frame}


\begin{frame}{Problem 2.8}
    Four students are deciding whether to buy a smoothie. Each student is willing to pay the following maximum amount for one smoothie:

\begin{itemize}
    \item Priya: \$7
    \item Marcus: \$5
    \item Elena: \$4
    \item Noah: \$2
\end{itemize}

The market supply is perfectly elastic at a price of \$3. 

Answer the following:

\begin{enumerate}
    \item Draw a clearly labeled diagram of the supply and demand for smoothies.
    \item Identify the equilibrium price and the number of smoothies sold. Clearly mark this equilibrium point on your graph.
    \item Calculate the total consumer surplus at the equilibrium price. Indicate the surplus areas on your graph.
    \item What is the total producer surplus in this market? 
\end{enumerate}
\end{frame}



\begin{frame}{Problem 2.8 Solution}
    
    \begin{center}
        \begin{figure}[h]
            \centering
            \includegraphics[width=0.7\linewidth]{Problem Book/q8.png}
        \end{figure}
    \end{center}

CS is $(7-3)+(5-3)+(4-3)=4+2+1=\$7.$ The PS is $(3-3)\times(3-0)=\$0.$
\end{frame}
%-------------------------------------------------------------


%--------------------------------------------------------------
\begin{frame}{Review: Elasticity}
\textbf{Example:} Movie Theater Ticket Prices vs. Movie Release Timing

Recall the graphical intution behind elasticity. Flatter means elastic and steeper means inelastic. Special cases include a vertical line, perfectly inelastic, and a horozontal line, perfectly elastic.

The Movie Theatre:
\begin{itemize}
    \item Charge higher prices on opening weekend
    \item Offer discounts on weekdays or lower prices weeks later
\end{itemize}

The Consumer:
\begin{itemize}
    \item Can stream the movie
    \item Can wait, pirate, or not watch the movie
\end{itemize}

What can we infer about the consumer's elasticity?
\end{frame}

\begin{frame}{Review: Elasticity}
\begin{itemize}
    \item \textbf{Opening weekend: Demand Inelastic}
    \item Fans strongly want to see movie now
    \item Fans less sensitive to increases in price
    \item \textbf{Later weeks: Demand More Elastic}
    \item More alternatives available such as streaming or hearing a summary from a friend
    \item Fans more sensitive to changes in price
\end{itemize}


\end{frame}

%--------------------------------------------------------------
\begin{frame}{Problem 2.22}
    In the market for dark chocolate, the income elasticity of demand is +2, and consumer incomes increase by 10\%. Which of the following statements is most accurate?

\begin{itemize}
    \item[a)] Demand for chocolate will increase by more than 10\%.
    \item[b)] Demand for chocolate will increase by exactly 10\%.
    \item[c)] Demand for chocolate will increase by less than 10\%.
    \item[d)] The price elasticity of demand must be greater than 1.
    \item[e)] There is not enough information to determine the change in quantity demanded.
\end{itemize}
\end{frame}

\begin{frame}{Problem 2.22 Solution}
    We are given:

\begin{itemize}
  \item Income elasticity of demand: \( \varepsilon_I = +2 \)
  \item Consumer incomes increase by 10\%
\end{itemize}

We apply the income elasticity formula:
\[
\varepsilon_I = \frac{\%\ \text{change in quantity demanded}}{\%\ \text{change in income}}
\]

Solving for the percentage change in quantity:
\[
\frac{\Delta Q}{Q} = \varepsilon_I \times \frac{\Delta I}{I} = 2 \times 10\% = 20\%
\]

\textbf{Interpretation:}  
Since the quantity demanded increases by 20\%, which is more than 10\%, the correct answer is:

\[
\boxed{\text{(a) Demand for chocolate will increase by more than 10\%.}}
\]

\end{frame}

%--------------------------------------------------------------


%--------------------------------------------------------------

\begin{frame}{Problem 2.29}
    You are given the following information:

\begin{itemize}
\item Price elasticity of demand for cigarettes at current prices is -0.5.
\item Current price of cigarettes is \$0.05 per cigarette.
\item Cigarettes are being purchased at a rate of 10 million per year.
\end{itemize}


Find a linear demand that fits this information, and graph
that demand curve.
\end{frame}

\begin{frame}{Problem 2.29 Solution}
     We want to find a linear demand function of the form:
\[
Q(P) = a - bP
\]

\textbf{Step 1: Use the elasticity formula}
The point elasticity of demand is:
\[
\varepsilon = \frac{\Delta Q}{\Delta P} \cdot \frac{P}{Q}
\]
Since $Q = a - bP$, we have $\frac{\Delta Q}{\Delta P} = -b$. Substituting into the elasticity formula:
\[
-0.5 = (-b) \cdot \frac{0.05}{10}
\Rightarrow -0.5 = -b \cdot 0.005
\Rightarrow b = \frac{0.5}{0.005} = 100
\]

\textbf{Step 2: Use the known point to solve for $a$}
We know $Q = 10$ when $P = 0.05$:
\[
10 = a - 100 \cdot 0.05 \Rightarrow 10 = a - 5 \Rightarrow a = 15
\]

\end{frame}

\begin{frame}{Problem 2.29 Solution}
   \textbf{Resulting Demand Function}
\[
Q(P) = 15 - 100P
\]
where $Q$ is measured in millions of cigarettes per year.

\begin{center}
    \begin{figure}[h]
        \centering
        \includegraphics[width=0.5\linewidth]{q29}
    \end{figure}
\end{center}

\end{frame}

\section{Consumer Theory}
\begin{frame}{Review: The Budget Line}
    \textbf{Example: College Meal Plan}

    A college student has \$200 per month for food. The student can choose to spend that money at the dining hall, grocery store, or a combination of both.

    How does an increase in dining hall prices, holding grocery store prices constant, effect the student's budget line?


\begin{itemize}
    \item The budget line: $p_1x_1 +p_2x_2 = I$
    \item The slope:  $-\frac{p_1}{p_2}$
    \item So if one price changes, the slope changes, but the line still passes through one fixed intercept.
    \item Therefore, the change in price results in a rotation of the budget line, but \textbf{not a shift} in the budget line
    \item The student's choices may change
\end{itemize}

\end{frame}



\begin{frame}{Problem 3.1}
  It is known that Mr. Hicks' income is \$500. It is also known that the prices are:

\[
P_1 = 
\begin{cases}
1 & \text{if } x_1 \leq 250 \\
2 & \text{if } x_1 > 250
\end{cases}
,\quad P_2 = 2
\]

Draw the consumer’s budget constraint.  
\end{frame}

\begin{frame}{Problem 3.1 Solution}
The budget line for $x_1\leq 250$: $1\cdot x_1+2x_2\leq500.$ Thus, the $x_2$ intercept is $\frac{I}{p_2}=\frac{500}{2}=250$ units. The slope of that part of the budget line in absolute value is $\frac{p_1}{p_2}=\frac{1}{2}.$ This is true until $x_1=250.$ Then, the remaining consumption of good 1 beyond 250 units is priced at \$2. At the point where $x_1=250$, since we only have an available income of $\$500-\$1\cdot 250=\$250$ to spend on good 2 we can only purchase at the same time 125 units of good 2. If we now want to increase consumption of good 1, the opportunity cost in terms of good 2 is no longer $\frac{1}{2}$ units of good 2, but $\frac{p_1}{p_2}=\frac{2}{2}=1$ units of good 2. So forgoing those 125 units, will allow us to buy exactly an additional (to the already 250 units we own) amount of 125 units of good 1. In total, we can buy \textbf{at most} 375 units of good 1.
\end{frame}

\begin{frame}{Problem 3.1 Solution}
    \begin{center}
    \begin{figure}[h]
        \centering
        \includegraphics[width=\linewidth]{Problem Book/Q1c3.png}
    \end{figure}
\end{center}
\end{frame}

%------------------------------------------------------------
\begin{frame}{Problem 3.4}
Dana is a consumer who enjoys two goods: loaves of bread (good 1) and cups of tea (good 2). Her total income is \$12. The prices of the goods are given by the vector $p = (2, 1)$, meaning a loaf of bread costs \$2 and a cup of tea costs \$1.

Dana shops at a local neighborhood supermarket, which currently offers a special promotion on bread: for every two loaves she purchases, she receives a third loaf for free. That is, for every 2 units of good 1 she buys, she gets 1 additional unit at no extra cost.

Draw Dana’s budget set. Be sure to account for the effect of the promotion when representing her affordable consumption bundle.
\end{frame}

\begin{frame}{Problem 3.4 Solution}
\begin{center}
    \begin{figure}[h]
        \centering
        \includegraphics[width=0.8\linewidth]{Problem Book/Q4c3.png}
     
    \end{figure}
\end{center}

\end{frame}

\begin{frame}{Review: Preferences and Properties}

    \textbf{Example: Choosing a Laptop}

    Andrew spent 12 hours training machine learning models this weekend and was horrified to learn that his friend ran the same models in an hour using a new laptop. Andrew is now going to buy a new laptop, we will simplify his comparison of laptops to CPU cores and battery life.

    How can we map Andrew's preferences of laptops to the properties of Completeness, Transitivity, Monotonicity, and Convexity?

\end{frame}

\begin{frame}{Review: Preferences and Properties}

\begin{itemize}
    \item Completeness: Andrew can rank any two laptops
    \item Transitivity: If a Macbook $>$ a PC and a PC $>$ a Chromebook, then a Macbook $>$ a Chromebook
    \item Monotonicity: More CPU cores and more battery is better
    \item Convexity: Andrew prefers a balanced laptop over an extreme (infinite cores and no battery is not optimal)
    \item Note, relations can be strict and weak. Andrew may be indifferent between two PCs and the completeness property holds
    \item We explain this phenomena using a \textbf{utility function}
\end{itemize}
\end{frame}

\begin{frame}{Problem 3.24}
    True/False/Not enough Information: A consumer with a utility function that has indifference curves bowed outwards as below, is maximizing utility at the bundle $c$.
\begin{figure}[H]
    \centering
    \includegraphics[width=0.5\linewidth]{Problem Book/indiff_curve1.jpg}
    \label{fig:indiff-curve}
\end{figure}
\end{frame}

\begin{frame}{Problem 3.24 Solution}
    False. At this current point $MRS<\frac{p_1}{p_2}$. In order to maximize, the consumer needs to reduce consumption of good 1 and increase consumption of good 2 so that $MRS=\frac{p_1}{p_2}$ which we know is the optimality condition since preferences are monotonic and strictly convex.
\end{frame}

\begin{frame}{Problem 3.19}
    Evaluate the following statements. Indicate whether each statement is \textbf{true or false, or not enough information} and \textbf{justify your answer}.

\begin{enumerate}
    \item A consumer cannot be indifferent between a bundle that lies \textbf{above their budget line} and a bundle that lies \textbf{below their budget line}.
    
    \item Two \textbf{indifference curves} of a given consumer \textbf{can never intersect}.
    
    \item If the prices of both goods $X$ and $Y$ drop by \textbf{50\%}, and the consumer originally spent their entire income on these two goods, then the consumer will \textbf{increase the quantity purchased} of \textbf{each good}.
\end{enumerate}

\end{frame}

\begin{frame}{Problem 3.19 Solution}
    \begin{enumerate}
    \item False. Affordability has nothing to do with preferences. You can have the IC intersect the budget line such that some bundles are strictly affordable and some are not. \pause 
    \item Not Enough Information. They cannot intersect, if preferences are transitive. If preferences are non-transitive, they might end up intersecting. (True would also be accepted with the right explanation.) \pause 
    \item Not Enough Information (False with right explanation is also a good answer). The budget set doubles in every direction, so the original bundle is still affordable. The substitution effect is zero. Thus, what happens to the quantities of both goods completely depends on the income effect which will be different if the goods are normal (both), or one normal one neutral, or one normal and one inferior. 
\end{enumerate}
\end{frame}

\end{document}